\documentclass[11pt,a4paper]{article}
\usepackage[margin=1in]{geometry}
\usepackage[T1]{fontenc}
\usepackage[utf8]{inputenc}
\usepackage{lmodern}
\usepackage{microtype}
\usepackage{enumitem}
\usepackage{xcolor}
\usepackage{hyperref}
\hypersetup{colorlinks=true, linkcolor=black, urlcolor=blue}
\setlength{\parindent}{0pt}
\setlength{\parskip}{0.55em}
\setlist[itemize]{leftmargin=1.5em,itemsep=0.2em,topsep=0.2em}
\setlist[enumerate]{leftmargin=1.7em,itemsep=0.2em,topsep=0.2em}
\pagenumbering{arabic}
\begin{document}
\begin{center}
{\LARGE\bfseries Company Website and Employee Portal Proposal}
\end{center}
\vspace{0.5em}

Hello,\par

My name is Fabricio Santana. I am a software engineer and technical lead with 20+ years of experience in backend development, software architecture, database integration, automated testing, cloud computing, and delivery of critical systems. I have also led large engineering teams and worked on projects where data consistency, business rules, validation, security, and maintainability are essential.\par

Your project is a strong match for my background because the main challenge is not only rebuilding pages with a new brand. The important part is separating the public website from the secure employee portal, preserving existing content and functions, and delivering a bilingual, maintainable system that your team can update without depending on developers for routine content changes.\par

My relevant approach for this kind of project is to combine public web delivery, secure authentication, internal workflow design, content management, bilingual editing strategy, deployment support, and technical handoff into one coherent delivery plan. This is important here because the public site and the employee portal have different user expectations, risk profiles, and maintenance needs.\par

\section*{Working methodology}

My delivery method is designed to reduce risk before development starts. The work begins with a fixed Discovery Sprint and then moves into a delivery package only after scope, priorities, responsibilities, dependencies, and acceptance criteria are clear.\par

The first step is a one-week Discovery Sprint with a fixed price of USD 500. In this phase, I can hold up to five remote meetings to understand the business problem, review the current platform, analyze the current data sources and workflows, map requirements, identify risks, review integrations, and define priorities. The deliverable is a practical work plan for the implementation, including recommended scope, technical approach, milestones, success criteria, and the most appropriate delivery package.\par

After discovery, the project can move into one of three fixed-price, four-week delivery packages. The package is selected according to scope, complexity, delivery speed, integration needs, quality requirements, and the amount of engineering support required. Some roles may participate part-time depending on the agreed scope, especially Scrum Master, technical leadership, test automation, DevOps, observability, and deployment support.\par

\textbf{Delivery package options}\par

\begin{itemize}
\item \textbf{Essential Delivery} --- USD 2,000, 4 weeks. Includes a Scrum Master and one full-stack engineer. This package is designed for a focused module, MVP, backend feature, import flow, or small integration with limited uncertainty. It is suitable when the scope is well defined, the number of integrations is limited, and the main goal is to deliver a reliable working feature without unnecessary overhead.
\item \textbf{Product Acceleration} --- USD 4,000, 4 weeks. Includes a Scrum Master, two full-stack engineers, and one engineer focused on test automation, DevOps, observability, and deployment. This package is recommended when faster delivery is needed or when the scope includes a product, API, integration, internal workflow, multiple features, or more demanding quality needs.
\item \textbf{Scale \& Intelligence} --- USD 6,000, 4 weeks. Includes a Technical lead, Scrum Master, two full-stack engineers, and one engineer focused on test automation, DevOps, observability, and deployment. This package is designed for complex systems, advanced data workflows, automation, auditability, security, observability, broader platform evolution, or stronger technical scale planning.
\end{itemize}

Role responsibilities are clear: the Scrum Master organizes scope, communication, ceremonies, progress tracking, and delivery alignment; full-stack engineers implement backend, frontend, API, and database changes as needed; the test automation, DevOps, observability, and deployment engineer supports automated tests, CI/CD, monitoring, release quality, and deployment; and the Technical lead, included in the Scale \& Intelligence package, owns architecture, technical decisions, code quality, and risk management for more complex projects.\par

For a project like this, the implementation setup can include public website architecture, portal authentication, content migration, bilingual content management, CMS or admin configuration, deployment support, and handoff work within the agreed package scope.\par

Each delivery cycle is executed with Scrum-inspired practices, frequent alignment, technical review, automated testing when applicable, and a demonstration of what was delivered. The client keeps ownership of the produced code, and the scope is agreed before execution to avoid open-ended work and unclear expectations.\par

\textbf{Construction defect warranty.} After delivery and acceptance, I include a 30-day warranty for construction defects related to the approved scope. This covers fixes for implementation errors, broken agreed behavior, or defects introduced by the delivered code. It does not cover new features, changes in business rules, new portal workflows, new content not included in the approved migration scope, third-party service changes, infrastructure changes outside the project, or requests that expand the original acceptance criteria.\par

\section*{Opportunity analysis}

This project should be treated as a website-and-portal architecture problem, not as a simple redesign. The public site and the employee portal serve different users and require different controls. The public site needs brand implementation, multilingual content, good editing ergonomics, reusable page structures, and preservation of existing functionality. The employee portal needs authentication, secure access, internal communication or resource-sharing workflows, and a maintenance model that does not create operational friction after launch.\par

The main technical risks are underestimating the portal scope, mixing internal authenticated workflows with public-site content in an unmanaged architecture, choosing an editing model that is hard to maintain in English and Spanish, and migrating current content and functions without regressions. The security model also needs to be explicit enough to cover the right level of authentication, password recovery, roles or permissions, session handling, and protected portal routes. For this reason, I recommend starting with a Discovery Sprint to separate the information architecture, portal feature boundary, content workflow, security model, and first operational release scope.\par

A practical first implementation cycle can start with the new bilingual public site, the agreed authenticated portal foundation, initial internal resource-sharing workflows, content migration for the defined scope, and the technical base required for future expansion. The exact first-release boundary should be confirmed during discovery so the project does not imply every possible future portal workflow in the first four-week cycle.\par

\section*{Proposed work plan}

I would approach the project in structured phases. The Discovery Sprint is the immediate recommended commitment, and implementation can move directly into a fixed delivery cycle once the public-site scope, portal scope, and acceptance criteria are confirmed.\par

\textbf{1. Discovery, information architecture, and technical design}\par

\begin{itemize}
\item Review the current website, preserved content and functionality, updated brand requirements, employee-portal expectations, and target release boundary.
\item Define the boundary between public-site features and portal-only features.
\item Define the preferred content-management model for English and Spanish updates, including who can edit, review, and publish content.
\item Produce a concrete implementation plan with architecture recommendation, migration scope, CMS or admin approach, security model, acceptance criteria, responsibilities, dependencies, risks, and package recommendation.
\end{itemize}

\textbf{2. Public website foundation and content structure}\par

\begin{itemize}
\item Design and implement the public website foundation according to the updated brand identity.
\item Define the page structure, navigation, bilingual content model, reusable sections, and maintainable editing workflow.
\item Preserve or rebuild the agreed current functionality.
\item Establish the technical base for deployment, SEO essentials, and long-term maintainability.
\end{itemize}

\textbf{3. Employee portal foundation and secure access}\par

\begin{itemize}
\item Implement the authentication and secure-access foundation for the employee portal.
\item Define the initial portal structure for internal communication and resource sharing according to the agreed scope.
\item Add role or permission handling, protected routes, password recovery, and session behavior according to the agreed first-version security model.
\item Keep the first implementation cycle focused on a bounded internal workflow rather than broad uncontrolled portal expansion.
\end{itemize}

\textbf{4. Content migration, bilingual workflow, and user experience refinement}\par

\begin{itemize}
\item Migrate the agreed current content and preserve key business functionality.
\item Implement the bilingual workflow for English and Spanish updates, including the agreed CMS/admin fields or editable page sections.
\item Refine the public and internal user experience according to the approved design direction.
\item Support the client's editing autonomy through documentation, practical handoff, and a content model designed for routine updates.
\end{itemize}

\textbf{5. Testing, deployment support, documentation, and handoff}\par

\begin{itemize}
\item Test the public site, portal authentication, key workflows, content model, and edge cases.
\item Add deployment support and operational visibility where appropriate.
\item Document the architecture, editing workflow, portal limitations, security assumptions, and maintenance steps.
\item Provide a handoff session or written notes so your team can manage content and operate the solution safely after launch.
\end{itemize}

\section*{Estimated timeline}

For the immediate next step, my recommendation is a 1-week Discovery Sprint. After discovery, the first implementation cycle can move into a fixed 4-week delivery package focused on the agreed first operational release. For a well-bounded bilingual public site plus a moderate authenticated employee portal, the likely recommendation is one Product Acceleration cycle. If the portal requires stronger roles and permissions, auditability, identity-provider integration, or broader internal workflows, the likely recommendation is one Scale \& Intelligence cycle.\par

A practical schedule would be:\par

\begin{itemize}
\item Week 1: Discovery Sprint, architecture definition, release boundary, bilingual content model, portal scope, acceptance criteria, and implementation plan.
\item Weeks 2 to 5: first fixed implementation cycle, covering the agreed first operational website and portal release.
\end{itemize}

If discovery shows that the first version is narrower than currently implied, the implementation may fit a smaller delivery shape. If it shows broader portal requirements, stronger security controls, more complex content migration, or additional integrations, I will recommend the larger package or additional cycles accordingly.\par

\section*{Availability and communication}

I can work remotely and keep communication in English through Workana messages, scheduled calls, and written progress updates. I usually prefer short alignment checkpoints during the week, especially after reviewing the current site and after each major implementation milestone.\par

This helps reduce misunderstandings and keeps the project aligned with your actual brand, content, portal workflow, and operational constraints.\par

\section*{Budget}

Based on the current understanding, my recommended immediate starting budget is USD 500 for the one-week Discovery Sprint. This gives you a concrete plan before committing to the full implementation budget, while keeping a clear path to delivery in the following implementation cycle.\par

After discovery, the first operational release will likely fit one of these packages:\par

\begin{itemize}
\item \textbf{Product Acceleration} at USD 4,000 for 4 weeks if the scope is a bounded bilingual public site plus a moderate authenticated employee portal with a practical CMS or admin editing workflow.
\item \textbf{Scale \& Intelligence} at USD 6,000 for 4 weeks if the scope includes stronger access control, auditability, identity-provider integration, broader internal workflow complexity, or more demanding content governance.
\end{itemize}

This assumes that the client provides current-site access, updated brand assets or design guidance, the bilingual content to preserve or migrate, and timely decisions about the portal scope, editing workflow, and hosting model. It also assumes the first release is focused on the agreed public-site pages and the agreed employee-portal foundation rather than every possible future internal feature.\par

The Discovery Sprint will refine the scope and confirm the correct implementation package. If the review shows a smaller and tightly bounded release scope, I can recommend a narrower delivery shape. If it shows broader portal complexity, more demanding security needs, or additional integrations, I will recommend the larger package or additional cycles.\par

\section*{Questions before final confirmation}

Before confirming the final scope, I would like to clarify a few points:\par

\begin{enumerate}
\item What content and functionality from the current website must be preserved exactly?
\item Is there already a new brand guide, design system, or page mockups?
\item What features must exist in the first version of the employee portal?
\item Will the portal need roles and permissions, or only basic authenticated access?
\item How should internal communication work in the first version: announcements, documents, resource directories, or another model?
\item What is the preferred editing workflow for English and Spanish content, and who should be able to edit, review, and publish changes?
\item Should the portal and the public site use the same CMS or admin system, or separate tools?
\item Is there an existing user directory or identity provider for employee authentication, or should the first version include standalone login and password recovery?
\item What hosting, deployment, and ownership model do you expect after launch?
\item Is there already an approved budget range and target launch date for the first operational version?
\end{enumerate}

Once these points are clear, I can confirm the final delivery plan and milestones.\par

\end{document}
